\documentclass[10pt]{article}
\usepackage[utf8]{inputenc}
\usepackage[T1]{fontenc}
\usepackage{amsmath}
\usepackage{amsfonts}
\usepackage{amssymb}
\usepackage[version=4]{mhchem}
\usepackage{stmaryrd}
\usepackage{bbold}

\begin{document}
Jum. - [1] С к о р и я к о в J. А., Дедекиндовы структуры с дополнениями и регулярные кольца, М., 1961; [2] Итоги науки. Алгебра. Топология. Геометрия, 1968 , М., 1970; [3] Ф о ф а н ов а Т. С., в сб.: Упорядочснные множества и решетки, в. 3, [Саратов], 1975 , c. $28-40 ;$ [4] M ur ra y F., N e u ma n n., "Ann. Math.", 1936 , v. 37, № 1, p. 116 - 229; [5] L o o m i s L. H., "Mem. Amer. Math. Soc.", 1955, № 18, p. 1-36; [6] M a e d a S., "J. Sci. Hiroshima Univ., Ser. A.", $1955, \mathrm{v}, 19, \mathcal{N o}_{2}$, p. $211-37$; [7] K a p l a n s k y I., "Ann. Math.", 1955, v. 61, № 3, p. 524-
541. OРТОНОРМИРОВАННАЯ СИСТЕМА -1) О. с. в е кт о р о в - множество $\left\{x_{\alpha}\right\}$ ненулевых векторов евклидова (гильбертова) пространства со скалярным произведением $(\cdot, \cdot)$ такое, что $\left(x_{\alpha}, x_{\beta}\right)=0$ при $\alpha \neq \beta$ (ортогональность) п $\left(x_{\alpha}, x_{\alpha}\right)=1$ (нормируемость).

\begin{enumerate}
  \setcounter{enumi}{1}
  \item 
  \begin{enumerate}
    \setcounter{enumii}{1}
    \item О. с. ф у нк ц и й - система $\left\{\varphi_{i}(x)\right\}$ функций пространства $L^{2}(X, S, \mu)$, являющаяся одновременно ортогональной и нормированной в $L^{2}(X, S, \mu)$, то есть
  \end{enumerate}
\end{enumerate}

\[
\int_{X} \varphi_{i}(x) \bar{\varphi}_{j}(x) d \mu= \begin{cases}0 & \text { при } i \neq j, \\ 1 \text { при } i=j\end{cases}
\]

(см. Нормированная система, ортогональная сисmема). В математич. литературе часто термин «ортогональная система» означает «ортонормированная система». При исследовании данной ортогональной системы ее нормированность не играет существенной роли. Тем не менее нормированность систем дает возможность более яснойі формулировки нек-рых теорем о сходимости рядов

\[
\sum_{k=1}^{\infty} c_{k} \varphi_{k}(x)
\]

в терминах поведения коэффициентов $\left\{c_{k}\right\}$. Такой теоремой является, напр., т е о р е м а Р ис с а- Ф иш е р: ряд

\[
\sum_{k=1}^{\infty} c_{k} \varphi_{k}(x)
\]

по ортонормированной в $L^{2}[a, b]$ системе $\left\{\varphi_{k}(x)\right\}_{k=1}^{\infty}$ сходится в метрике пространства $L^{2}[a, b]$ тогда и только тогда, когда

\[
\sum_{k=1}^{\infty}\left|c_{k}\right|^{2}<\infty
\]

Лum.: [1] К о л м о г о р о в А. Н., Фо м и н С. В., Элементы теории функций и функционального анализа, 5 изд., М., $1981 ;[2]$ К а ч м а ж С., Ш т е й н г а у з $\Gamma$., Теория ортогональных рядов, пер. с нем., М., 1958. А. А. Талалян.

OPTОЦЕНТР т р у у го л ь и к а - точка пересечения трех высот треугольника. О. треугольника лежит на Эйлера прямой. Середины трех сторон, середины отрезков, соединяющих О. с тремя вершинами, и основания высот треугольника лежат на одной окружности. О. является центром окружности, вписанной в о р тоце т т и ч е с кий т р у голы и ик, т. е. треугольник, вершинами к-рого являются основания высот данного.

П. С. Моденов.

ОСВЕШҢЕИЯ ЗАДАЧА - задача определения минимального числа направлений пучков параллельных лучеї или числа источников, освещающих всю границу выпуклого тела. Пусть $K$ - выпуклое тело $n$-мерного линейного пространства $\mathbb{R}^{n}$, a bd $K$ и int $K$ - cooтветственно граница и внутренность его, причем bd $K \neq$ $\neq \varnothing$. Наиболее известны следуюцие О. 3.

\begin{enumerate}
  \item Пусть $l$ - нек-рое направление в пространстве $\mathbb{R}^{n}$. Точка $x \in \mathrm{bd} K$ наз. о с в е щ е н о й и з в н н а п р а в л е н и е м $l$, если прямая, проходящая через $x$ параллельно $l$, проходит через нек-рую точку $y \in \operatorname{int} K$ и направление вектора $\overrightarrow{x y}$ совпадает с $l$. Ищется минимальное число $c(K)$ направлений в пространстве $\mathbb{R}^{n}$, достаточное для освецения в этом смысле всего множества bd $K$.

  \item Пусть $z$ - нек-рая точка множества $\mathbb{R}^{n} \backslash K$. Точка $x \in \mathrm{bd} K$ наз. о с в е е н н о й и з в н т т ч ч- к о ӥ $z$, если прямая, определяемая точками $z$ п $x$, проходит через нек-рую точку $y \in$ int $K$ п векторы $\overrightarrow{x y}$ и $\overrightarrow{z y}$ одинаково направлены. Ипется минимальное число $c^{\prime}(K)$ точек из $\mathbb{R}^{n} \backslash K$, достаточное для освещения в этом смысле всего множества bd $K$.

  \item Пусть $z$ - нек-рая точка множества bd $K$. Точка $x \in \mathrm{bd} K$ наз. о с в е і е н н о й и з н у т р и т о чк о й $z \neq x$, если прямая, определяемая точками $z$ и $x$, проходит через нек-рую точку $y \in$ int $K$ и векторы $\overrightarrow{x y}$ и $\overrightarrow{z y}$ противоположно направлены. Ищется минимальное число $p(K)$ точек из $\mathrm{bd} K$, достаточное для освещения изнутри всего множества bd $K$.

  \item Система точек $Z=\{z \mid z \in \mathrm{bd} K\}$ наз. ф и к с и р ов ан н о й для $K$, если она обладает свойствами: a) $Z$ достаточна для освещения изнутри всего множества bd $K$; б) $Z$ не обладает никаким собственным подмножеством, достаточным для освецения изнутри множества bd $K$. Ищется максимальное число $p^{\prime}(K)$ точек фиксированной системы для тела $K \subset \mathbb{R}^{n}$.

\end{enumerate}

Задача 1) была поставлена в связи с Xадвигера әипотезой (см. [1]): минимальное число тел $b(K)$, гомотетичных ограниченному $K$ с коэфффициентом гомотетии $k, 0<k<1$, достаточное для покрытия $K$, удовлетворяет неравенству $n+1<b(K) \leqslant 2^{n}$, причем значение $b(K)=2^{n}$ характеризует параллелепипед. Для ограниченного $K \subset \mathbb{R}^{n} c(K)=b(K)$. Если $K$ неограниченно, то $c(K) \leqslant b(K)$ и существуют такие тела, что $c(K)<b(K)$ или $c(K)=b(K)=\infty$ (см. [1]).

Задача 2) поставлена в связи с задачей 1). Для ограниченного $K \subset \mathbb{R}^{n}$ верно равенство $c(K)=c^{\prime}(K)$. Если же $K$ неограниченно, то $c^{\prime}(K) \leqslant b(K)$ и $c(K) \leqslant$ $\leqslant c^{\prime}(K)$. Число $c^{\prime}(K)$ для любого неограниченного $K \subset \mathbb{R}^{3}$ принимает одно из значений: $1,2,3,4, \infty$ (cM. [1]).

Решение задачи 3) имеет вид: число $p(K)$ определено тогда и только тогда, когда $K$ отлично от конуса. В этом случае

\[
2 \leqslant p(K) \leqslant n+1,
\]

причем $p(K)=n+1$ характеризует $n$-мерный симплекс пространства $\mathbb{R}^{n}$ (см. [1]).

Для задачи 4) (см. [2]) предполагается, что при ограниченном $K \subset \mathbb{R}^{n}$ верно неравенство

\[
p^{\prime}(K) \leqslant 2^{n} \text {. }
\]

Каждая из О. з. тесно связана с нек-рым специальным покрытием тела $K$ (см. [1]).

Лum.: [1] Б о лтян с ки й В. Г., С о лт ан П. С., Комбинаторная геометрия различных классов выпуклых множеств, Kиш., 1978; [2] G r ü n b a u m B., "Acta math. Acad. sci
hung.», 1964, v. 15, p. 161-63. OCEBOЙ
ВEKTOP, а к с и а л в н ы й В е к т о р, п с е в д о в е к т о р, вектор в ориентированном пространстве, к-рый при изменении ориентации пространства на противоположную преобразуется в противоположный вектор. Пример О. в.- векторное про$\begin{array}{ll}\text { изведение векторов. } & \boldsymbol{Б} C \vartheta-3\end{array}$

ОСНАЩЕННОЕ ГИЛЬБЕРТОВО ПРОСТРАНСТВО гильбертово пространство $H$ с выделенным в нем линейным всюду плотным подмножеством $\mathrm{D} \subset H$, на к-ром задана структура топологического векторного пространства так, что вложение непрерывно. Это вложение порождает непрерывное вложение сопряженных пространств $\Phi^{\prime} \subset H^{\prime}$ и цепочку непрерывных вложений $\Phi \subset H \subset \Phi^{\prime}$ (с помоцью стандартного отождествления $\left.H^{\prime}=H\right)$. Наиболее содержательным является случай, когда оснащение Ф - ядерное nросmранство. Здесь верно следующее усиление спектральной теоремы для самосопряженных операторов, дей-


\end{document}